\chapter{Experimental Setup}\label{chap:setup}

This chapter fixes the models, benchmarks, sparse-attention baselines,
and speed-measurement protocol used throughout
Chapter~\ref{chap:results} and Chapter~\ref{chap:analysis}.


\section{Models}\label{sec:models}

All experiments use two open-weight, decoder-only language models at
the $\approx$8B-parameter scale: Llama-3.1-8B-Instruct
\cite{grattafiori2024llama3} and Qwen3-8B~\cite{yang2025qwen3}.
Both use rotary positional embedding (RoPE) and grouped-query
attention (GQA); Qwen3-8B additionally applies QK-Norm to the query
and key projections.  All experiments run in \texttt{bfloat16} on a
single NVIDIA RTX A6000 (48~GiB) unless noted otherwise.


\section{Benchmarks}\label{sec:benchmarks}

DCT-Page is evaluated on three long-context benchmark suites.  All
prompt-formatting follows the chat-template conventions baked into
each model's tokenizer, with the exceptions documented in
Section~\ref{sec:bench-longbench}.

\subsection{RULER}\label{sec:bench-ruler}

RULER~\cite{hsieh2024ruler} is a synthetic long-context test suite of
13 tasks, used here at a default context length of 32K tokens.
The tasks fall into four families: needle-in-a-haystack retrieval
(NIAH Single 1--3, NIAH MultiKey 1--3, NIAH MultiValue, NIAH
MultiQuery), variable tracking (VT), word-frequency extraction
(Common Words Extraction, Frequent Words Extraction), and
long-document question answering (QA-1, QA-2).
Each task is scored by its native exact-match or substring metric.

\subsection{LongBench v1 and v2}\label{sec:bench-longbench}

\paragraph{LongBench v1}~\cite{bai2024longbench} provides a suite of
English long-context tasks across single-document QA, multi-document
QA, summarization, and few-shot in-context learning.  We evaluate on
the six tasks most representative of these categories:
NarrativeQA, Qasper, MultiFieldQA-en, 2WikiMultihopQA, GovReport, and
TriviaQA.  Metrics are the official per-task ones (F1 for QA, ROUGE-L
for summarization, classification accuracy for in-context learning).
TriviaQA is evaluated without a chat template; the other five tasks
use the model's chat format.

\paragraph{LongBench v2}~\cite{bai2025longbenchv2} is a 503-question
multiple-choice benchmark with four answer options per question.  We
report overall accuracy and break it down by difficulty
(\texttt{easy} / \texttt{hard}) and by context length
(\texttt{short} / \texttt{medium} / \texttt{long}).


\section{Baselines and Hyperparameters}\label{sec:baselines}

DCT-Page is compared against five sparse-attention methods organized
into two paradigms.  All five leave prefill unchanged and modify only
decode-time attention; their algorithmic details and provenance are
covered in Chapter~\ref{chap:background}.  This section fixes the
configuration with which each method is run in
Chapter~\ref{chap:results}.

\paragraph{Page-based baselines.}
Quest~\cite{tang2024quest},
SeerAttention-R~\cite{gao2025seerattention}, and
InfLLM~\cite{xiao2024infllm} share the fixed-page paradigm with
DCT-Page.  For each we match the \emph{selected-token budget}
$B \cdot P = 64 \cdot 32 = 2{,}048$ tokens per decode step, the same
budget DCT-Page uses in Section~\ref{sec:config}; settings within
that budget follow each method's upstream recommendation.  Table
\ref{tab:baseline-config} summarizes the per-method configuration.

\paragraph{Cluster-based reference.}
Multipole Attention~\cite{hooper2025multipole} organizes keys into a
$k$-means cluster tree rather than fixed-size pages, and its per-step
work cannot be matched against the page-based budget in a one-to-one
fashion.  Multipole typically achieves the highest accuracy across
our benchmarks but at substantially longer decode latency, so we
exclude it from the main accuracy tables and instead devote
\S\ref{sec:results-multipole} to its accuracy--throughput trade-off
against DCT-Page.  Multipole is run with a single-level
configuration of $3.125\%$ leaf-cluster retention --- conceptually
matched to DCT-Page's $P{=}32$ page granularity --- and a per-level
token-budget threshold of $2{,}048$.  Clusters are formed at the end
of prefill, and additional clusters are added during decode as the
cache grows.

\paragraph{Full-KV reference.}
Every accuracy table additionally includes a Full-KV row, which is
each model run with its stock dense decode attention.  This is the
accuracy upper bound that all sparse methods aim to approach at
lower decode-time cost.

\paragraph{Model coverage.}
Page-based baselines are run on whichever model the upstream code
supports: SeerAttention-R is reported on Qwen3-8B only (the released
gates are trained on the Qwen family, with no equivalent
Llama-3 checkpoint), InfLLM on Llama-3.1-8B-Instruct only (the
upstream targets the Llama family), and Quest on both models.
Multipole and Full-KV are reported on both models.

\begin{table}[ht]
  \centering
  \caption[Per-method configuration used in the experiments.]
  {Per-method configuration used in the experiments.  ``Budget''
   denotes the per-decode-step token budget that the method's
   selection rule must satisfy.}
  \label{tab:baseline-config}
  \small
  \begin{tabular}{l p{0.55\textwidth} l}
    \toprule
    Method            & Key settings                                      & Models \\
    \midrule
    DCT-Page          & $P{=}32,\ B{=}64,\ C/P{=}1/8,\ \rho{=}\gamma{=}\max$ & both \\
    Quest             & page size $32$, budget $2{,}048$ tokens           & both \\
    SeerAttention-R   & token-budget mode, budget $2{,}048$ tokens        & Qwen3 only \\
    InfLLM            & block size $32$, top-$64$ blocks, sink $32$,
                        local $4{,}096$, repr.\ top-$4$                   & Llama only \\
    Multipole         & $3.125\%$ leaf retention ($\equiv P{=}32$),
                        threshold $2{,}048$, centroid replacement         & both \\
    Full-KV           & dense decode attention                            & both \\
    \bottomrule
  \end{tabular}
\end{table}


\section{Speed-Measurement Protocol}\label{sec:speed-protocol}

All speed numbers are collected on a single NVIDIA RTX A6000 (48~GiB)
from one unified profiler that drives the model with random
token inputs and records both the end-to-end decode latency and a
per-stage breakdown in a single run.  Two execution paths are
compared:

\begin{itemize}
  \item \emph{Full-KV} --- the model with its stock dense decode
        attention, served by upstream FlashInfer's paged-attention
        primitive.
  \item \emph{DCT-Page} --- the production configuration introduced
        in Section~\ref{sec:kernels}, in which DCT-Page's
        scoring/selection kernels feed per-(batch, kv-head) page
        indices into the same upstream FlashInfer primitive via a
        virtual-batching scheme that requires no custom fork of the
        attention runtime.
\end{itemize}

\paragraph{CUDA graphs.}
Every speed measurement reported in this thesis uses CUDA-graph
capture for the decode loop.  This pins the indices/indptr buffers
that the paged-attention primitive reads and removes per-step Python
and indices-buffer allocation overheads that would otherwise dominate
at short context length.

\paragraph{Throughput measurement.}
For each $(L, \text{batch})$ cell with $L \in \{8, 16, 32, 64, 128\}\text{K}$
tokens and batch sizes $\{1, 2\}$, the profiler runs $8$ warm-up
replays followed by $128$ timed decode-step replays of the captured
CUDA graph, bracketed by a pair of CUDA events that span the entire
timed loop.  Per-step wall-clock latency is the elapsed time of this
bracket divided by the number of replays, and is reported as two
quantities: (i)~\emph{attention-only} latency (the time spent inside
the attention sub-module, including scoring/selection for DCT-Page),
and (ii)~\emph{end-to-end} per-step latency (the same chain extended
through the MLP, layer-norm, and projection).
Speedups in Section~\ref{sec:results-speed} are reported as the ratio
of the Full-KV latency to the DCT-Page latency at the same
$(L, \text{batch})$.

\paragraph{Per-stage breakdown.}
Recording CUDA events between substeps inside the captured graph is
not supported on the torch + CUDA build used in this thesis (event
records issued during graph capture either error out or fail to
produce timings on replay), so the per-stage breakdown of
Section~\ref{sec:results-speed} is collected through the kernel-level
path instead.  Concretely, the same timed replay loop is wrapped in a
\texttt{torch.profiler} session that captures CUPTI kernel activity;
each kernel launched inside one replay window is then bucketed by
its kernel name into one of the eight decode substeps --- QKV
projection, RoPE + cache append, page-region view (a stride-only
operation, no kernel), proxy update, page scoring, fused
top-$k$/KV-assembly, the upstream FlashInfer attention kernel, and
the output projection.  The bucketing dictionary is anchored to
specific kernel symbols emitted by FlashInfer, the DCT-Page Triton
kernels, and the upstream paged-attention primitive; a residual
``all-kernel reconciliation'' check ensures that the sum of bucketed
times agrees with the total CUPTI activity to within sub-millisecond
tolerance, so unbucketed kernels (typically a renamed cuBLASLt
symbol) cannot silently disappear from the table.

For Full-KV the same bucketing collapses to four populated substeps
(QKV, RoPE + cache append, FlashInfer attention, output projection);
all other DCT-Page substeps are structurally absent.  In both modes,
QKV projection and the output projection operate on the single query
token and are identical across the two paths, so the reported
attention-time budget is the sum of the remaining substeps:
\begin{equation*}
  t_\text{Full-KV}^{\text{attn}}
  \;\;\text{vs.}\;\;
  t_\text{DCT-Page}^{\text{proxy}}
  + t_\text{DCT-Page}^{\text{score}}
  + t_\text{DCT-Page}^{\text{topk+assemble}}
  + t_\text{DCT-Page}^{\text{attn}}.
\end{equation*}
Reporting the comparison as a sum keeps the overhead of DCT-Page's
bookkeeping kernels charged against its measured speedup.
